\subsection{Définition des systèmes d'ouverture/parachutes}
\label{subsec:def_para}

Dans l'objectif de récupération de nos deux étages, nous devons définir et valider nos deux systèmes de sauvegarde des lanceurs. Ceux-ci sont définis par des trappes latérales actionnées par un servomoteur chacun.\\


\begin{table}[h!]
\centering
\begin{tabular}{|p{1.2cm}|p{4cm}|p{10cm}|}
\hline
Règles & Type 2 & Détails du cahier des charges \\ \hline
REC1  & Système récupération & Chaque étage doit être équipé d’un système ralentisseur fiable permettant de réduire sa vitesse de descente. L’éjection du/des ralentisseurs doit être franche. \\ \hline
REC2  & Système récupération & Les systèmes ralentisseurs des étages et de tout autre élément éjecté doivent permettre une arrivée au sol à une vitesse verticale comprise entre 5 m/s et 15 m/s, y compris dans le cas dégradé où les deux étages ne se sont pas séparés. \\ \hline
REC3  & Cohérence du système & Les instants de déploiement des systèmes ralentisseurs doivent être compatibles avec l’expérience menée par le club. \\ \hline
REC4  & Séparation transversale & Non utilisé. \\ \hline
REC5  & Séparation latérale & La case contenant le système de récupération doit rester opérationnelle lorsqu’elle supporte en compression longitudinale une force : $F = 2 \times \text{Accélération}_{\text{Max}} \times M_{\text{sup}}$ (numériquement, l’accélération en m/s$^2$ et la masse en kg donnent $F$ en newtons). \\ \hline
REC6  & Assemblage trappe & En position fermée, la porte latérale ne doit pas dépasser du profil de l’étage. \\ \hline
REC7  & Tenue en torsion & La porte ne doit pas s’ouvrir ou se bloquer lorsqu’on applique un couple de torsion de 1 N·m entre le haut et le bas de l’étage. \\ \hline
REC8  & Vibrations & L’accélération et les vibrations de la fusée ne doivent pas modifier le fonctionnement des systèmes de récupération. \\ \hline
REC9  & Tenue mécanique globale & Les systèmes de récupération utilisés doivent auparavant avoir été testés avec succès. \\ \hline
REC10 & Solidité du parachute & Les ralentisseurs doivent être suffisamment solides pour résister au choc à l’ouverture. \\ \hline
REC11 & Anneau anti-torche & Dans le cas de l’utilisation d’un parachute, celui-ci doit être équipé d’un anneau anti-torche (également appelé glisseur). \\ \hline
\end{tabular}
\end{table}



La contrainte imposée par le cahier des charges est une vitesse de redescente comprise entre 5 et 15m/s pour les trois configurations (étage 1, étage 2 et non séparé).
De par les dimensions des deux étages, le premier sera en charge du retour sous parachute de l'ensemble non séparé.
\begin{equation}
\begin{aligned}
    V_d &= \sqrt{\frac{2 M g}{\rho_0 C_x S}} \\
    %V_d &= \sqrt{\frac{2 \times 7.2 \times 9.81}{1.3 \times 1 \times 0.65}}
\end{aligned}
\end{equation}

Avec la règle REC2, nous devons dimmensionner les deux parachutes pour les cas suivants :

\begin{table}[h!]
\centering
\begin{tabular}{|c|c|c|c|}
\hline
Cas d'étude & Pas de séparation & Étage 1 & Étage 2 \\ \hline
Masse de l'ensemble & 8,080kg & 4,759kg & 3,321kg \\ \hline
Dimensions du parachute utilisé & a=b=360mm & a=b=360mm & a=b=310mm \\ \hline
Vitesse sous parachute & 13,8 m/s & 9,6 m/s & 10,7 m/s\\ \hline
\end{tabular}
\end{table}
